% qkit-syllabus LaTeX preamble.
% Loaded via the YAML include-in-header field in index.qmd.
%
% Everything that does NOT need a value out of the YAML header lives here.
% The two pieces that do -- the optional logo and the course-information
% table -- live in before-body.tex instead, because include-in-header is
% copied in verbatim and Pandoc never interpolates $variables$ into it.

\usepackage{titlesec}
\usepackage{titling}
\usepackage{etoolbox}
\usepackage{fancyhdr}
\usepackage{lastpage}
\usepackage{enumitem}
\usepackage{booktabs}
\usepackage{array}
% tabularx builds the course-information table in before-body.tex; array
% supplies the >{...} column modifiers that make its two columns ragged.
\usepackage{tabularx}

% ---------------------------------------------------------------------
% Typewriter and sans families.
%
% `fontfamily: mathpazo` sets the ROMAN family to Palatino and leaves the
% other two at Computer Modern. Under the T1 encoding Pandoc selects,
% those resolve to the METAFONT ec fonts, which pdflatex embeds as Type 3
% bitmaps -- every `code span` in the PDF then renders fuzzy and is not
% searchable. Latin Modern is the Type 1 equivalent and ships with
% TinyTeX. Redefining the two defaults is the fix rather than loading
% lmodern, because Pandoc loads mathpazo before this file and a later
% \usepackage{lmodern} would take the roman family back off Palatino.
%
% Check with `pdffonts index.pdf`: every row should read Type 1.
% ---------------------------------------------------------------------
\renewcommand{\ttdefault}{lmtt}
\renewcommand{\sfdefault}{lmss}

% ---------------------------------------------------------------------
% Title block.
%
% titling supplies \pretitle / \posttitle, which before-body.tex uses to
% hang an optional logo above the title. It also replaces article's
% \maketitle, and that is the load-bearing part: article's \@maketitle
% opens with \newpage, so a logo emitted before \maketitle would be
% stranded on a page of its own. titling's version has no \newpage, so
% logo, title and information table share the first page.
% ---------------------------------------------------------------------
\setlength{\droptitle}{-.25cm}

% ---------------------------------------------------------------------
% Headings.
%
% Sections are unnumbered (number-sections: false in the YAML), so the
% spacing below is all that separates one part of the syllabus from the
% next -- keep the "before" skip clearly larger than the "after" one or
% headings stop reading as breaks.
% ---------------------------------------------------------------------
\titlespacing\section{0pt}{12pt plus 4pt minus 2pt}{6pt plus 2pt minus 2pt}
\titlespacing\subsection{0pt}{12pt plus 4pt minus 2pt}{6pt plus 2pt minus 2pt}
\titleformat*{\subsubsection}{\normalsize\itshape}

% ---------------------------------------------------------------------
% Running head and footer.
%
% \@title and \@date rather than $title$ and $date$: Quarto's title.tex
% partial issues \title{} and \date{} AFTER this file is included, so the
% macros are undefined at the point \rhead is read. That does not matter,
% because \rhead only stores tokens -- both are expanded at shipout, by
% which time Quarto has defined them. Keeping them here rather than in
% before-body.tex leaves the whole page-style setup in one file.
%
% \@date is whatever before-body.tex last passed to \date{} -- the `term:`
% key when it is set, otherwise Quarto's own `date:`. Delete both and the
% head ends in a stranded dash.
%
% The one caveat is \thanks{}: Quarto appends it to \@title, so a syllabus
% that sets `thanks:` in the YAML drags that footnote into the running
% head. Set the running head by hand with \rhead{} in that case.
% ---------------------------------------------------------------------
\pagestyle{fancy}
\fancyhf{}
\renewcommand{\headrulewidth}{0.3pt}
\renewcommand{\footrulewidth}{0pt}
\makeatletter
\rhead{\footnotesize \@title{} --- \@date}
\makeatother
\cfoot{\small \thepage/\pageref*{LastPage}}

% Page 1 already carries the title and the information table, so it takes
% the page number but neither the rule nor the running head.
% before-body.tex applies this with \thispagestyle.
\fancypagestyle{firststyle}{%
  \fancyhf{}%
  \renewcommand{\headrulewidth}{0pt}%
  \fancyfoot[C]{\small \thepage/\pageref*{LastPage}}%
}

% ---------------------------------------------------------------------
% Reading list.
%
% Citations go through natbib rather than citeproc (see the YAML header),
% so the heading over the generated list is \refname. Quarto's
% `reference-section-title` key has no effect on the natbib path -- change
% the line below instead.
% ---------------------------------------------------------------------
\renewcommand{\refname}{Recommended Readings}

% Drop natbib's comma before the year, giving "(McKinney 2022)" rather
% than "(McKinney, 2022)".
\setcitestyle{aysep={}}

% Give the generated entries a little air.
\AtBeginEnvironment{thebibliography}{%
  \setlength{\itemsep}{4pt}%
}
